% Generated draft: requires mathematical and source review.
\documentclass[11pt]{article}
\usepackage[margin=25mm]{geometry}
\usepackage{fontspec}
\setmainfont{Latin Modern Roman}
\usepackage{amsmath,amssymb,mathtools,unicode-math}
\setmathfont{Latin Modern Math}
\usepackage{enumitem,needspace,xurl,hyperref}
\setlength{\parindent}{0pt}
\setlength{\parskip}{5pt}
\setlength{\emergencystretch}{3em}
\newcommand{\N}{\mathbb{N}}
\newcommand{\Z}{\mathbb{Z}}
\newcommand{\Q}{\mathbb{Q}}
\newcommand{\R}{\mathbb{R}}
\newcommand{\C}{\mathbb{C}}
\newcommand{\F}{\mathbb{F}}
\newcommand{\A}{\mathbb{A}}
\newcommand{\PP}{\mathbb{P}}
\newcommand{\E}{\mathbb{E}}
\newcommand{\eps}{\varepsilon}
\newcommand{\dd}{\,\mathrm{d}}
\begin{document}
\section*{UnsolvedMath ID: 30004344}
\textbf{Source problem code:} catalogue entry
\subsection{Definitions and mathematical statement}
The supplied record retains a conditional theorem extract, rather than a recoverable standalone open problem. Its supplied verification note identifies the extract as a theorem proved assuming the Bloch--Beilinson conjecture. The following records its mathematical content schematically; the missing hypotheses and definitions prevent a complete theorem statement.

For each integer \(n\) with \(1\leq n\leq m\), write \(A_n\) for the unspecified Chow coefficient group containing the special-cycle classes \([Z(T,\varphi)]\). This notation is only a placeholder: the underlying variety, codimension, coefficient ring, and dependence on the level or other datum \(S\) are not supplied. The recorded generating series is
\[
\phi_n(\tau,\varphi,S)
 =\sum_{T\geq 0}[Z(T,\varphi)]q^T.
\]
Here \(\tau\), \(\varphi\), \(S\), \(m\), the index set for \(T\), the meaning of \(T\geq0\), and the Fourier monomial \(q^T\) refer to source data that are absent from the supplied extract. In particular, the relevant number field, matrix sizes, positivity conditions at its embeddings, and analytic domain are not specified. No definition of the classes for boundary or degenerate indices is provided.

The retained assertion is that, assuming the intended Bloch--Beilinson conjecture, \(\phi_n\) is a Chow-valued Hilbert--Siegel modular form for every \(1\leq n\leq m\). More explicitly, its stated scalar-valued consequence is that for every linear functional \(\ell\) from the relevant Chow coefficient space to \(\C\),
\[
\ell\bigl(\phi_n(\tau,\varphi,S)\bigr)
 :=\sum_{T\geq0}\ell\bigl([Z(T,\varphi)]\bigr)q^T
\]
is the Fourier expansion of a complex Hilbert--Siegel modular form. This quantifier retains the source wording; the scalar field for linearity and any necessary scalar extension of the Chow group are unspecified. The modular weight, group, level, representation or character, and holomorphy and boundary requirements are also absent.

The extract additionally asserts a pullback formula, but supplies neither the formula nor its maps and hypotheses; that assertion cannot be stated precisely from this record. The literature assessment refers instead to a Beilinson--Bloch injectivity conjecture. Its relationship to the Bloch--Beilinson hypothesis in the theorem extract is not specified and is not assumed here. Thus these expressions preserve the conditional assertion without supplying a complete modularity theorem or an additional open conjecture.
\subsection{Short English statement}
The retained extract says that, under an unspecified Bloch--Beilinson hypothesis, special Chow-cycle generating series are Hilbert--Siegel modular for every \(1\leq n\leq m\), including after applying every linear functional to their coefficients. It also mentions an unavailable pullback formula. Essential geometric data and modular-form parameters are missing. The supplied verification identifies this as a conditional theorem extract, while the registry labels it partially solved; a precise standalone open problem cannot be recovered.
\subsection{Sources}
\begin{itemize}
\item Supplied local UnsolvedMath registry record, permanent ID \(30004344\), title ``Conditional Modularity Theorem for Chow Cycles,'' source code \texttt{catalogue entry}, in the Oberwolfach Reports Open Problems collection. The record gives proposed year 2019 and names no proposer. Its statement status is \texttt{unrecoverable}; its supplied problem-status label is \texttt{partially\_solved}.
\item Supplied report citation: \emph{Modular Forms (2020)}, Oberwolfach Reports, DOI 10.4171/owr/2019/57, pp.\ 3529--3606. \url{https://doi.org/10.4171/owr/2019/57}. The supplied verification note says that the official report presents the extraction as a theorem proved assuming the Bloch--Beilinson conjecture, rather than a separately posed open problem. No separate clean statement is supplied.
\item Supplied literature source: ``Maeda, modularity of special cycles.'' \url{https://doi.org/10.4153/S000843952000020X}. The supplied record does not provide a full bibliographic citation or theorem number.
\item Supplied literature source: ``OWR report.'' \url{https://doi.org/10.4171/owr/2019/57}.
\item Literature assessment dated exactly \texttt{2026-08-22}: the supplied curation describes Chow-valued modularity as proved assuming the Beilinson--Bloch injectivity conjecture, with unconditional results in important special cases, and says that the general conditional hypothesis remains unresolved. This is supplied dated curation, not an independently verified assessment. The background explicitly requests independent verification of time-sensitive status and cited claims.
\end{itemize}
\end{document}
